\subsection{The physical affine lifting identity}

Use component order $(r,\phi,z)$ and the full matrix space $\mathscr A$.
Let $\chi_f(s)=\widehat\chi_f(\nu s)$ equal one on $s\le1/\nu$
and zero on $s\ge1.1/\nu$. For $A\in\mathscr A$ set
\[
 A^y=\chi_f A,\qquad B=(A^y)'+\mu A^y/s,\qquad
 M(A)=\sup_{0<s\le1.1/\nu}s^{-5/2}
                \sum_{j=0}^6|(s\partial_s)^jA^y(s)|.
\]
The cutoff derivative is included in $B$; the matrices are not
continued beyond their support. All assertions are uniform for
$\mu\in K\Subset(0,1/4)$ and the fixed cutoff.

Write $a=1+\nu s$, $x_j=(q-jc,\eta)$, $c=\nu^{-2}$, and
\[
 V_j=a^{-2}e^{-|x_j|^2/(4a)}Jx_j,\quad
 p_j=-a^{-3}e^{-|x_j|^2/(2a)},\quad
 \psi_j=-2a^{-1}e^{-|x_j|^2/(4a)},\qquad J(q,\eta)=(-\eta,q).
\]
Embed $V_j=J\nabla\psi_j$ with zero azimuthal component and define
\[
 \begin{split}
 f_A&=-\sum_j[(B_{\rm mer}x_j)\cdot\nabla V_j+BV_j],\\
 e_A&=\sum_j[2BV_j+2\nu A^y:D^2V_j
                         -(A^y+(A^y)^T)\nabla p_j],\\
 h_A^p&=\sum_j[V_j,A^y_{\rm mer}x_j],\\
 \psi_A^p&=\sum_j[({\rm tr}A^y_{\rm mer})\psi_j
                     -(A^y_{\rm mer}x_j)\cdot\nabla\psi_j].
 \end{split}
\]
The full matrix, including its hoop entry, is used before restriction
to the $(r,z)$-plane; symmetry of $A$ is not assumed. The identity
$A^y_{\rm mer}J=J[({\rm tr}A^y_{\rm mer})I-(A^y_{\rm mer})^T]$
gives $h_A^p=J\nabla\psi_A^p$.
The sums and their required derivatives converge by Gaussian decay.
With $\mathcal K=\eta\partial_\eta+2s\partial_s$, put
\[
 \widehat h_A=\sqrt s\,[(h_A^p)_r e_r+
       ((h_A^p)_z+\epsilon(1-\mathcal K)\psi_A^p)e_z],
 \qquad F_A=s^{3/2}f_A,\qquad E_A=s^{3/2}e_A.
\]
The physical cylindrical streamfunction is $r\psi_A^p$, with
$r=R/\sqrt s$, so $\widehat h_A$ is exactly divergence-free, vanishes
near the axis, and has young support and finite energy. The forces
are unprojected; the physical body inverse supplies the Cartesian
Leray projection. All these maps depend linearly on $A$.

\begin{lemma}\label{ret:physical-affine}
There is a real physical body force $\mathfrak r_A$, linear in $A$,
such that on the exact background $U$,
\[
 (L_{{\rm f},\mu}+\mathbb P\operatorname{div}M_U)\widehat h_A
       =\mathbb P(F_A+E_A+\mathfrak r_A).
\]
For $0<\epsilon,\nu\le1$, with $\delta_U=\|U-U_{\rm raw}\|_{W^{1,\infty}}$,
\begin{equation}\label{ret:general-affine-bounds}
 \begin{split}
 R^{-1}(\|F_A\|_{Y_2}+\|E_A\|_{Y_2}
       +\|\widehat h_A\|_{Y_2}+\|\nabla\widehat h_A\|_{Y_2})
       &\le P_\nu M(A),\\
 R^{-1}\|\mathfrak r_A\|_{Y_2}
       &\le P_\nu(\epsilon+\delta_{\rm base}+\delta_U)M(A).
 \end{split}
\end{equation}
Here $\delta_{\rm base}$ is the fixed envelope for the train-to-Gaussian
and cross-center norms on $0\le s\le1.1/\nu$, with the required
finite spatial and conormal derivative budget.
\end{lemma}
\begin{proof}
Write $v$ for the planar train in this calculation. The Cartesian
bracket calculation, with the full matrix retained, gives
\[
 (\partial_s+\mu/s-\nu\Delta)h_A^p
 +(v\cdot\nabla)h_A^p+(h_A^p\cdot\nabla)v
       =f_A+e_A+\nabla\Psi_A+r_{\rm per},\qquad
 \Psi_A=\sum_j(A^y_{\rm mer}x_j)\cdot\nabla p_j.
\]
Here $r_{\rm per}$ consists of distinct-center convection products
and $(v-\sum_jV_j)$ terms. Gaussian summation and the fixed-band
train comparison bound its unnormalized cell norms, with the finite
derivative budget, by $C\delta_{\rm base}M(A)s^{5/2}$.

The exact ring-chart rules are
\[
 \partial_r=\sqrt s(\partial_q-\epsilon\mathcal K),\qquad
 \partial_z=\sqrt s\partial_\eta,\qquad
 (\partial_\tau-y\cdot\nabla/2-1/2)(\sqrt s b)
       =s^{3/2}\partial_s b.
\]
The physical Laplacian of $\sqrt s b$ is $s^{3/2}$ times
\[
 [\Delta_{q,\eta}-2\epsilon(\mathcal K+1)\partial_q
       +\epsilon^2(\mathcal K+1)^2]b-\epsilon^2b_r e_r.
\]
Together with the stream completion, these identities put every
geometric difference into $\epsilon M(A)$ times a fixed Gaussian
polynomial of age size at most $C(s^{3/2}+s^{5/2})$ in age-force units.
The pressure is completed before projection:
\[
 s^{3/2}\nabla_{q,\eta}\Psi_A
 =\nabla_y(s\Psi_A)+\epsilon s^{3/2}(\mathcal K+2)\Psi_Ae_r.
\]
The first term is a global gradient with adequate decay; the second
belongs to the geometric error. Thus the physical pressure tail
remains accounted for without identifying the two Leray projections.
The remaining convection difference is exactly
\[
 (U-U_{\rm raw})\cdot\nabla\widehat h_A
            +\widehat h_A\cdot\nabla(U-U_{\rm raw}),
\]
which is bounded by $\delta_U$ times the asserted energy norm.

The Gaussian cell norms and fixed spatial derivatives are bounded by
$CM(A)s^{3/2}$ for $f_A$ and by $CM(A)s^{5/2}$ for both
$h_A^p$ and $\psi_A^p$.
The residual $e_A$ has size $CM(A)(s^{3/2}+s^{5/2})$.
For body fields the physical conversion is
\[
 R^{-2}\|s^{3/2}g\|_{Y_2}^2
   =\pi\int_0^{1.1/\nu}s^{1/2}
                   \|g(s)\|_{L^2(dq/c\,d\eta)}^2ds.
\]
Thus the relevant age powers, including $s^{7/2}$ and $s^{11/2}$,
are integrable at zero; their upper-end integrals and the normalization
$c^{-1/2}=\nu$ cost only $P_\nu$. The same calculation with the
physical spatial Jacobian $s^{-5/2}$ bounds the lift and its gradient.
The smooth cutoff introduces no boundary source.

For the derivative budget, the bracket uses $A^y$, the stream
completion uses one conormal derivative, and one physical spatial
derivative uses two. Its Laplacian uses three; one further spatial
derivative of the residual uses four. The declared six suffice.
These are body and spatial estimates; no bound for arbitrary pure
phase derivatives is asserted. The maps are continuous in $\mu\in K$.
\end{proof}

There are two uses of this single estimate. For generic $A\in\mathscr A$,
$M(A)\le C e^{1.1a_s/\nu}\|A\|_{\mathscr A}$. For the actual jet
$A=\mathcal LT$, the separated-kernel conormal bound gives instead
$M(A)\le P_\nu\epsilon^3N_1(\chi_oT_1)$, since $\chi_f A=A$.
Together with $\delta_U\le P_\nu\rho$ from the background bounds,
this gives the actual-range remainder without the generic exponential
envelope. For $a_\mu(s)=s^{5/2}\operatorname{diag}(12,-3,-9)/(\mu+5/2)$,
the formulas give the unit affine force used in the signed return.

The unprojected force also satisfies the stronger weighted estimate:
\[
 R^{-1}\|g^{-1}F_A\|_{Y_2}\le P_\nu\|A\|_{\mathscr A},\qquad
 g(s)=\exp[\gamma s/(1+\nu s)].
\]
Indeed $a_s<\gamma/16$ and $s\le1.1/\nu$ give
$g(s)^{-1}e^{a_ss}\le e^{-(\gamma/3-a_s)s}$; inserting this
pointwise bound in the same Gaussian integrals proves the estimate.
Multiplication by $g^{-1}$ is not commuted through the Cartesian
Leray projection; the near comparison keeps that projection explicit.
